\chapter{Observation Equations and Partition Algebra}
\label{ch:observation_equations}

\papernote{Source paper: \emph{Observational Partition Algebra} --- \texttt{docs/observation-equations/}}

\begin{chapterabstract}
Measurement, physical process, and observation are mathematically identical within the partition framework. For the partition equation $\Gamma_1 \oplus P(\omega) \to \Gamma_2$, the output is simultaneously: (i)~a physical result, (ii)~a measurable observable, and (iii)~a computational output. The observational identity theorem is proven. Categorical distance is shown to be independent of physical (spatial) distance, enabling opacity-independent measurement. Enzymes are information catalysts that reduce categorical distance. The Cellular Partition Language (CPL) is specified. Disease states correspond to specific categorical coherence signatures, enabling diagnostic identification through oscillatory fingerprinting.
\end{chapterabstract}

\input{../docs/observation-equations/cellular-observation-equations.tex}
